% CSC5515 G14 final report. Every number below is computed from the released data and the outputs in
% outputs/ by the scripts in src/ (see results/ and docs/week7_data_report.md). Build: xelatex report.tex (twice).
\documentclass[10pt,twocolumn]{article}
\usepackage[a4paper,margin=1.9cm,columnsep=0.7cm]{geometry}
\usepackage{fontspec}
\usepackage{polyglossia}
\setdefaultlanguage{english}
\setotherlanguage{bengali}
\setmainfont{Latin Modern Roman}
\newfontfamily\bengalifont[Script=Bengali,Scale=0.92]{Noto Sans Bengali}
\usepackage{amsmath}
\usepackage{booktabs}
\usepackage{multirow}
\usepackage{xcolor}
\usepackage{caption}
\captionsetup{font=small,labelfont=bf}
\usepackage[hidelinks]{hyperref}
\newcommand{\bn}[1]{\textbengali{#1}}
\newcommand{\pmsd}[1]{{\scriptsize$\pm$#1}}
\setlength{\tabcolsep}{4pt}
\raggedbottom

\title{\Large\textbf{Small Fine-Tuned Models, Not Prompted LLMs: Translating Bangla Regional Dialects into
Standard Bangla, and What Dialect Distance Explains}}
\author{Esm-e Moula Chowdhury Abha (26-94089-2)\\[2pt]
\small CSC5515 Natural Language Processing, Group G14\\
\small\url{https://github.com/EsmeAbha/preserving-bangla-dialect}}
\date{}

\begin{document}
\maketitle

\begin{abstract}
Bangla regional dialects such as Chittagonian and Sylheti differ strongly from standard Bangla, and Bangla
language technology built on standard text fails on them. We study dialect-to-standard translation on
Vashantor (five dialects) with a leakage-free split, and on ONUBAD as an out-of-domain test, after finding that
323 ONUBAD pairs share their standard sentence with Vashantor training data. With CPU-only compute we compare
copying the input, a word-replacement dictionary, BanglaT5, prompted Qwen2.5-3B-Instruct (zero- and 5-shot,
three runs each on validation) and Qwen2.5-0.5B-Instruct fine-tuned with LoRA. Prompting the 3B model does not
beat copying the input («zs_sub» and «fs_sub» against «copy_sub» chrF on a 500-pair test subset), whereas the
0.5B model after «ft_short» reaches «ft_sub», «ft_sub_gain» chrF above 5-shot prompting. Neither LLM reaches
the dictionary («dict_test» chrF on the full test set) or BanglaT5 («bt5_test»), and BanglaT5 is the only
system that clearly beats the dictionary out of domain. Dialect distance, 100 minus the chrF of the copy
baseline, ranks the dialects by difficulty for every system, and per sentence it explains almost all of the
difficulty for prompting ($\rho=«rho_zs»$) but less for trained systems ($\rho=«rho_ft»$). A dictionary
learned on four dialects closes about half of the gap on a fifth, unseen dialect. An automatic error profile
shows that prompted models keep dialect words or invent new ones, whereas trained models replace dialect words
but the fine-tuned LLM often picks the wrong standard word.
\end{abstract}

\section{Introduction}
Bangla is spoken in many regional dialects. Chittagonian and Sylheti differ so much from standard Bangla
that many speakers from other regions cannot follow them: the same sentence uses different words, verb
endings and pronouns while its meaning stays the same (Table~\ref{tab:examples}). Language technology for
Bangla is built almost entirely on standard written Bangla, so dialect text in social-media posts, customer
messages or call-centre transcripts is often misread. A system that rewrites dialect text into standard Bangla
\emph{without changing its meaning} lets existing Bangla tools work on such text.

Instruction-tuned large language models (LLMs) are an obvious candidate, because they can be prompted with a
short instruction and a few examples. We test whether this is enough, and compare prompting with fine-tuning
a small open-weight LLM with low-rank adapters (LoRA) \cite{lora,qlora}, with a fine-tuned Bangla
sequence-to-sequence model (BanglaT5) \cite{banglat5}, and with two simple baselines that turn out to matter:
copying the input, and a word-replacement dictionary learned from the training pairs. We ask three questions:

\begin{itemize}\setlength\itemsep{1pt}
\item \textbf{RQ1.} Does a fine-tuned LLM preserve meaning better than zero-shot and few-shot prompting?
\item \textbf{RQ2.} Which dialects are hardest, and is the difficulty explained by the word-level distance of
the dialect from standard Bangla, measured as $100-\mathrm{chrF}$ of the copy baseline?
\item \textbf{RQ3.} How much of what is learned for four dialects transfers to a fifth, unseen dialect?
\end{itemize}

\noindent Our contributions are:
(i) a leakage-free re-split of Vashantor \cite{vashantor} that keeps all dialect versions of a standard
sentence in one split, and the finding that 323 of the 2,796 ONUBAD \cite{onubad} pairs, often used as a
separate test set, share their standard sentence with Vashantor training data;
(ii) evidence that prompting a 3B instruction-tuned LLM is no better than copying the input, whereas a 0.5B
LLM fine-tuned with LoRA on CPU improves on it by «ft_sub_gain» chrF;
(iii) a word-replacement dictionary baseline that is hard to beat («dict_test» chrF on test, ahead of both LLMs and level with BanglaT5), which earlier work did not report;
(iv) a dialect-distance measure that predicts per-sentence difficulty for every system (RQ2) and a
leave-one-dialect-out study (RQ3); and (v) an automatic word-level error profile that separates
\emph{kept dialect words} from \emph{invented words}.
All code, splits and outputs are public.

\begin{table}[t]\centering\small
\begin{tabular}{@{}l@{\hspace{6pt}}l@{}}\toprule
Standard & \bn{আজকে আমার মন ভালো নেই}\\\midrule
Chittagong & \bn{আযিয়া আর মন বালা নাই?}\\
Noakhali & \bn{আইজ্জা আর মন ভালা নাই}\\
Sylhet & \bn{আইজকু আমার মন ভালা নায়}\\
Barishal & \bn{আইজ মোর মনডা ভালোনা}\\
Mymensingh & \bn{আইজ আমার মন ভালা নাই}\\\bottomrule
\end{tabular}
\caption{One Vashantor standard sentence (``I am not in a good mood today'') and its five dialect versions.}\label{tab:examples}
\end{table}

\section{Related work}
Faria et al.\ \cite{vashantor} released Vashantor, a parallel corpus of five dialects and standard Bangla,
and fine-tuned mT5 and BanglaT5 \cite{banglat5}, reporting BLEU from 36.75 (Chittagong) to 69.06
(Mymensingh). Jawad et al.\ \cite{jawad} benchmarked zero- and few-shot prompting of LLMs on 600 YouTube
comments in four dialects and found that transliteration helps closed models. Sami et al.\ \cite{sami} used
retrieval-augmented prompting for the opposite direction, standard-to-dialect. Sultana et al.\ \cite{onubad}
collected ONUBAD, a separate corpus of Chittagong, Sylhet and Barishal sentences. We add a fine-tuned LLM,
simple but strong baselines, a significance-tested comparison over several runs, an explicit distance measure
for dialect difficulty, and leave-one-dialect-out transfer.

\section{Data}
\paragraph{Vashantor.} Mendeley Data V2 (doi:10.17632/bj5jgk878b.2, CC BY 4.0). Each pair joins a dialect
sentence with its standard sentence; Banglish, English and region columns are not used. The official split
has 1,875 / 250 / 375 pairs per dialect. Cleaning (Unicode NFC, zero-width characters removed, whitespace
collapsed) removed no pairs: the release has no empty and no duplicate pairs.

\paragraph{Leakage-free split.} The same standard sentences are translated into all five dialects, so a test
sentence must not reach training through another dialect. We group pairs by their standard sentence,
ignoring punctuation and spaces; a group found in more than one split is moved entirely to its test (else
validation) split. One sentence (\bn{আপনি কেমন আছেন?}) was in both training and validation; its five
training pairs were moved, giving 1,874 / 251 / 375 pairs per dialect (Table~\ref{tab:data}).
Seventeen groups have fewer than five members because the standard sentence is spelled differently in some
files (e.g.\ \bn{নিষেধ}/\bn{নিশেধ}); each still lies in one split.

\begin{table}[t]\centering\footnotesize
\begin{tabular}{@{}lrrrr@{}}\toprule
Split & Pairs & Per dialect & Std.\ sent. & Words \\\midrule
Train & 9,370 & 1,874 & 1,876 & 6.7 \\
Validation & 1,255 & 251 & 250 & 5.5 \\
Test & 1,875 & 375 & 381 & 7.3 \\\midrule
ONUBAD (all) & 2,796 & 930--935 & -- & -- \\
ONUBAD (novel) & 2,437 & -- & -- & -- \\\bottomrule
\end{tabular}
\caption{Data after cleaning and alignment (words: mean length of the standard sentence). ONUBAD covers
Chittagong, Sylhet and Barishal only.}\label{tab:data}
\end{table}

\paragraph{ONUBAD.} Mendeley Data V2 (doi:10.17632/6ft99kf89b.2, CC BY 4.0); we use only its sentence
file. After removing 144 duplicates there are 2,796 pairs. \textbf{323 pairs have a standard sentence that
also occurs in the Vashantor training set} and 126 a dialect sentence that does, so they are not really
out-of-domain. We report ONUBAD on the 2,437 \emph{novel} pairs without this overlap.

\section{Systems}
\paragraph{Copy.} The input is returned unchanged. Its chrF against the reference measures how far a
dialect is from standard Bangla; we define \emph{dialect distance} $=100-\mathrm{chrF}_{\text{copy}}$.

\paragraph{Dictionary.} Per dialect, the two token sequences of each training pair are aligned with a
longest-matching-block alignment; aligned one-to-one and one-to-many replacements are counted, and each dialect
word is replaced by its most frequent standard form at test time. Unknown words are copied. A \emph{pooled}
variant learns one table from all dialects.

\paragraph{Prompted LLM.} Qwen2.5-3B-Instruct \cite{qwen} in bf16, with its chat template and the
instruction ``\emph{Rewrite the following sentence from the \{dialect\} dialect into standard Bangla. Keep the
meaning unchanged. Output only the standard Bangla sentence.}'' followed by the sentence. Zero-shot uses three
paraphrases of the instruction (three runs); 5-shot adds five training pairs of the same dialect as earlier
chat turns, with three random draws (three runs).

\paragraph{Fine-tuned LLM.} Qwen2.5-0.5B-Instruct with LoRA adapters (rank 32, $\alpha=64$, dropout 0.05) on
all linear layers (17.6M trainable parameters, 3.4\%), trained on the same prompt as zero-shot with the loss
on the answer tokens only. AdamW, learning rate $2\cdot10^{-4}$, cosine schedule with 3\% warm-up,
effective batch 16, maximum length 256 tokens, bf16 autocast. We report «ft_epochs». Validation chrF on a fixed
stratified 250-pair subset is measured twice per epoch and the best adapter is kept. Training ran on four CPU
cores (about 22\,s per step).

\paragraph{BanglaT5.} \texttt{csebuetnlp/banglat5} \cite{banglat5}, fully fine-tuned in fp32 on inputs
``\texttt{<Dialect>: <sentence>}'' normalised with the BanglaT5 normaliser (learning rate $3\cdot10^{-4}$,
batch 16). On CPU we could afford one epoch (63 minutes).

\paragraph{Decoding.} Greedy, at most 128 new tokens, stopping at the end-of-turn token; the first non-empty
line is kept and prefixes such as ``Standard Bangla:'' are removed. Empty outputs are scored, not dropped.

\section{Evaluation}
chrF \cite{chrf} (sacreBLEU) is the primary metric, because dialect and standard forms often share
characters but not whole words; BLEU is reported for comparison with \cite{vashantor}; LaBSE \cite{labse}
cosine similarity between output and reference is a meaning measure that does not depend on spelling.
Systems with several runs are reported as mean $\pm$ standard deviation. Differences are tested with a
paired bootstrap over sentences (1,000 resamples, chrF recomputed from sufficient statistics, runs averaged)
and reported with 95\% confidence intervals (CI). Prompts, the 5-shot draw and the checkpoint were chosen on
validation only; the test split was scored once, at the end. Because the prompted 3B model needs about
1.4--5\,s per sentence on CPU, it was run on the test split on a fixed stratified subset of 500 pairs
(100 per dialect), using the best zero-shot instruction and the best 5-shot draw on validation; all other
systems are reported on the same subset as well as on the full test set.

\section{Results}
\subsection{Development on validation}
Prompting is stable but weak. Over three instruction paraphrases, zero-shot Qwen2.5-3B scores
39.9\pmsd{0.8} chrF on validation, and over three example draws 5-shot scores 41.0\pmsd{0.4}, against 41.0
for copying the input; the 5-shot gain over zero-shot is small but significant ($+1.1$, CI $[+0.6, +1.7]$).
We kept the best paraphrase (40.5) and the best draw (41.4) for the test runs. The fine-tuned 0.5B model
scores «ft_val» on validation. Dictionary and BanglaT5 score 71.0 and 69.9.

\begin{table*}[t]\centering\small
\begin{tabular}{@{}lrrrrrrrr@{}}\toprule
System & Chit. & Noak. & Sylh. & Bari. & Myme. & chrF & BLEU & LaBSE\\\midrule
Copy input & 32.0 & 39.1 & 52.6 & 40.0 & 62.3 & 45.4 & 12.3 & 0.71\\
\addlinespace[1pt]
Qwen2.5-3B zero-shot & 32.6 & 38.1 & 51.7 & 39.2 & 59.1 & 44.3 & 11.1 & 0.71\\
Qwen2.5-3B 5-shot & 34.7 & 39.2 & 50.2 & 42.0 & 56.9 & 44.8 & 12.8 & 0.73\\
Qwen2.5-0.5B + LoRA & 45.7 & 53.8 & 61.8 & 60.6 & 76.8 & 59.9 & 30.2 & 0.83\\
Dictionary (per dialect) & 55.4 & 60.5 & 69.8 & \textbf{77.5} & \textbf{89.1} & 70.7 & 45.3 & 0.91\\
BanglaT5, 1 epoch & \textbf{58.3} & \textbf{62.4} & \textbf{74.6} & 71.8 & 86.9 & \textbf{71.0} & \textbf{47.3} & \textbf{0.92}\\
\bottomrule\end{tabular}
\caption{RQ1 on the fixed 500-pair stratified test subset (100 per dialect): prompting versus fine-tuning. The prompted model uses the best instruction and the best 5-shot draw on validation.}\label{tab:sub500}
\end{table*}

\begin{table*}[t]\centering\small
\begin{tabular}{@{}lrrrrrrrr@{}}\toprule
System & Chit. & Noak. & Sylh. & Bari. & Myme. & chrF & BLEU & LaBSE\\\midrule
Copy input & 31.0 & 38.5 & 51.5 & 38.2 & 59.4 & 43.7 & 10.8 & 0.70\\
\addlinespace[1pt]
Dictionary (per dialect) & 56.1 & 60.8 & 68.5 & \textbf{78.2} & \textbf{87.9} & 70.3 & 44.4 & 0.90\\
Dictionary (pooled) & \textbf{58.3} & 63.5 & 70.9 & 78.2 & 87.0 & \textbf{71.6} & 46.0 & \textbf{0.91}\\
BanglaT5, 1 epoch & 57.2 & \textbf{63.8} & \textbf{75.0} & 73.0 & 85.3 & 70.9 & \textbf{46.9} & \textbf{0.91}\\
Qwen2.5-0.5B + LoRA & 44.5 & 54.6 & 61.7 & 59.4 & 76.9 & 59.4 & 30.7 & 0.83\\
\bottomrule\end{tabular}
\caption{Full Vashantor test set (1,875 pairs, 375 per dialect): chrF per dialect and overall, BLEU and LaBSE. Best value per column in bold.}\label{tab:main}
\end{table*}


\subsection{RQ1: prompting versus fine-tuning}
On the 500-pair test subset (Table~\ref{tab:sub500}, differences in Table~\ref{tab:sig}), zero-shot
Qwen2.5-3B scores «zs_sub» chrF, \emph{below} copying the input at «copy_sub» («zs_vs_copy»). Five examples
add little («fs_vs_zs») and do not lift it above copying («fs_vs_copy»); LaBSE agrees (0.71--0.73 against
0.71 for copying). Fine-tuning changes the picture: Qwen2.5-0.5B, six times smaller, reaches «ft_sub» chrF
after «ft_short», clearly above 5-shot («ft_vs_fs») and zero-shot prompting («ft_vs_zs»), with LaBSE
«ft_labse».
For models of this size the answer to RQ1 is clearly yes.

Fine-tuning a general LLM is, however, not the best use of the training data (Table~\ref{tab:main}). On the
full test set the 0.5B model («ft_test» chrF) is behind both the dictionary («ft_vs_dict») and BanglaT5
(«ft_vs_bt5»). BanglaT5 after one epoch and the dictionary are statistically tied («bt5_vs_dict»). One pooled
dictionary is better than per-dialect tables (per-dialect minus pooled: «dict_vs_pool»), although on
validation the two were equal ($+0.1$, CI $[-0.4, +0.6]$): many replacements are shared between dialects, so
dialect-specific tables only lose data.

\begin{table*}[t]\centering\small
\begin{tabular}{@{}llrcr@{}}\toprule
Set & Comparison & $\Delta$chrF & 95\% CI & $n$\\\midrule
Sub. & 3B 5-shot $-$ 3B zero-shot & $+0.45$ & [$-$0.53, 1.40] & 500\\
Sub. & 3B 5-shot $-$ Copy & $-0.63$ & [$-$1.71, 0.44] & 500\\
Sub. & 3B zero-shot $-$ Copy & $-1.08$ & [$-$1.81, $-$0.37] & 500\\
Sub. & Qwen-0.5B+LoRA $-$ 3B 5-shot & $+15.19$ & [13.51, 16.81] & 500\\
Sub. & Qwen-0.5B+LoRA $-$ 3B zero-shot & $+15.64$ & [14.16, 17.18] & 500\\
\addlinespace[2pt]
Test & Qwen-0.5B+LoRA $-$ Dictionary & $-10.89$ & [$-$11.68, $-$10.10] & 1875\\
Test & Qwen-0.5B+LoRA $-$ BanglaT5 & $-11.54$ & [$-$12.39, $-$10.63] & 1875\\
Test & BanglaT5 $-$ Dictionary & $+0.64$ & [$-$0.17, 1.45] & 1875\\
Test & Dictionary $-$ Dict.\ pooled & $-1.30$ & [$-$1.62, $-$0.94] & 1875\\
Test & Qwen-0.5B+LoRA $-$ Copy & $+15.69$ & [14.83, 16.45] & 1875\\
\addlinespace[2pt]
ONU. & Qwen-0.5B+LoRA $-$ Dictionary & $-12.40$ & [$-$13.53, $-$11.32] & 2437\\
ONU. & BanglaT5 $-$ Dictionary & $+2.94$ & [1.87, 4.04] & 2437\\
\bottomrule\end{tabular}
\caption{Paired bootstrap differences in chrF (1,000 resamples) on the 500-pair subset (Sub.), the full test set (Test) and novel ONUBAD (ONU.).}\label{tab:sig}
\end{table*}


\subsection{RQ2: distance and difficulty}
On test, dialect distance orders the dialects Chittagong (69.0), Barishal (61.8), Noakhali (61.5), Sylhet
(48.5), Mymensingh (40.6). On validation Sylhet was third (60.2), so the order of the middle dialects depends
on the sentences; Chittagong is always the farthest and Mymensingh the closest. Every system finds
Chittagong hardest and Mymensingh easiest, and over the five dialects distance and chrF correlate at
$-0.7$ to $-1.0$ for all systems (Table~\ref{tab:rq2}). Per sentence, distance explains almost all of the
difficulty for the prompted model ($\rho=«rho_zs»$ zero-shot, $«rho_fs»$ 5-shot): it changes little, so it
succeeds only where the input was already close to the reference. For trained systems the correlation is
weaker ($\rho=«rho_dict»$ dictionary, $«rho_bt5»$ BanglaT5, $«rho_ft»$ fine-tuned LLM): training removes part
of the dependence on distance, but far dialects stay hard. The answer to RQ2 is yes: copy-based distance is a
cheap and reliable predictor of difficulty, at the dialect and at the sentence level.

\begin{table}[t]\centering\footnotesize
\begin{tabular}{@{}lrcr@{}}\toprule
System & $\rho$ (sentences) & 95\% CI & $\rho$ (dialects)\\\midrule
\multicolumn{4}{@{}l}{\emph{Full test set} ($n=1{,}875$)}\\
Dictionary & $-0.58$ & $[-0.62, -0.55]$ & $-0.70$\\
Dictionary (pooled) & $-0.57$ & $[-0.61, -0.54]$ & $-0.70$\\
BanglaT5 & $-0.58$ & $[-0.61, -0.54]$ & $-0.90$\\
Qwen-0.5B + LoRA & $-0.55$ & $[-0.58, -0.51]$ & $-0.90$\\
\multicolumn{4}{@{}l}{\emph{500-pair subset}}\\
Qwen2.5-3B zero-shot & $-0.90$ & $[-0.92, -0.87]$ & $-1.00$\\
Qwen2.5-3B 5-shot & $-0.77$ & $[-0.81, -0.73]$ & $-1.00$\\
Qwen-0.5B + LoRA & $-0.57$ & $[-0.63, -0.51]$ & $-1.00$\\
\bottomrule\end{tabular}
\caption{RQ2: Spearman correlation between dialect distance ($100-$copy chrF) and chrF, per sentence (with bootstrap CI) and over the five dialect means. More negative = difficulty follows distance more closely.}\label{tab:rq2}
\end{table}


\subsection{RQ3: unseen dialects}
A pooled dictionary learned from four dialects and applied to the fifth closes 44--61\% of the gap between
copying and a dictionary that has seen the dialect (Table~\ref{tab:lodo}). The shared part of dialect
vocabulary is therefore large: \bn{ভালা} (standard \bn{ভালো}, good) occurs in the training data of four
dialects, and the pronouns \bn{মুই}, \bn{মোর} (I, my) in both Barishal and Sylhet. Barishal recovers least on
test (44\%) and Chittagong on validation (23\%), so which dialect transfers worst is not stable across
sentence samples. Leave-one-dialect-out fine-tuning of the
LLMs needs GPU time and is left for future work.

\begin{table}[t]\centering\small
\begin{tabular}{@{}lrrrr@{}}\toprule
Held-out & Copy & Unseen & Seen & Recovered\\\midrule
Chittagong & 31.0 & 44.6 & 58.3 & 50\%\\
Noakhali & 38.5 & 52.4 & 63.5 & 56\%\\
Sylhet & 51.5 & 63.4 & 70.9 & 61\%\\
Barishal & 38.2 & 55.8 & 78.2 & 44\%\\
Mymensingh & 59.4 & 74.6 & 87.0 & 55\%\\
\bottomrule\end{tabular}
\caption{RQ3 with the pooled dictionary on test (chrF on the held-out dialect). \emph{Unseen}: learned from the other four dialects; \emph{Seen}: learned from all five. Recovered = share of the gap between copy and seen that is closed without any data from the dialect.}\label{tab:lodo}
\end{table}


\subsection{Out of domain}
All systems lose on novel ONUBAD (Table~\ref{tab:onubad}): the dictionary «dict_drop_onu» chrF (to
«dict_onu»), BanglaT5 «bt5_drop_onu» (to «bt5_onu») and the fine-tuned LLM «ft_drop_onu» (to «ft_onu»), while
copying the input scores «copy_onu». ONUBAD contains only three dialects, including the hardest, so part of
the drop comes from the dialect mix. Out of domain BanglaT5 beats the dictionary («bt5_vs_dict_onu»): a model
of the language generalises to unseen words, a lookup table cannot.

\begin{table}[t]\centering\footnotesize
\begin{tabular}{@{}lrrrrrr@{}}\toprule
System & Chit. & Sylh. & Bari. & chrF & BLEU & LaBSE\\\midrule
Copy & 29.8 & 39.4 & 43.7 & 37.6 & 10.4 & 0.71\\
\addlinespace[1pt]
Dictionary & 50.0 & 62.9 & 65.6 & 59.5 & 31.5 & 0.85\\
BanglaT5 & \textbf{51.1} & \textbf{65.4} & \textbf{70.9} & \textbf{62.4} & \textbf{35.9} & \textbf{0.87}\\
Qwen-0.5B+LoRA & 38.0 & 49.9 & 53.8 & 47.1 & 17.4 & 0.76\\
\bottomrule\end{tabular}
\caption{Out of domain: the 2,437 novel ONUBAD pairs (no sentence shared with Vashantor training).}\label{tab:onubad}
\end{table}


\subsection{Error analysis}
Table~\ref{tab:errors} shows \emph{how} the systems fail. Zero-shot prompting keeps «zs_kept»\% of the
dialect words and produces only «zs_prod»\% of the standard words that were needed; examples move this to
«fs_kept»\% and «fs_prod»\%, but at the cost of more invented words («fs_novel»\% of the output). The
dictionary keeps «dict_kept»\% of the dialect words (unknown words are copied) but rarely invents
(«dict_novel»\%). BanglaT5 and the fine-tuned LLM both replace almost every dialect word («bt5_kept»\% and
«ft_kept»\% kept), but the LLM produces fewer of the needed words («ft_prod»\% against «bt5_prod»\%) and more
novel ones («ft_novel»\% against «bt5_novel»\%): it has learned \emph{that} a word must change more reliably
than \emph{what} it must change into. Negation disagrees with the reference in «bt5_neg»--«ft_neg»\% of the
trained systems' outputs. Table~\ref{tab:qual} gives examples.

\begin{table}[t]\centering\footnotesize
\begin{tabular}{@{}lrrrrr@{}}\toprule
System & Exact & Kept & Produced & Novel & Neg.\\\midrule
Copy & 1.0 & 100.0 & 0.0 & 0.0 & 7.0\\
Dictionary & 17.0 & 30.1 & 53.1 & 11.2 & 4.2\\
BanglaT5 & 18.0 & 9.4 & 56.5 & 20.8 & 3.8\\
3B zero-shot & 0.6 & 62.1 & 6.3 & 22.9 & 5.8\\
3B 5-shot & 1.0 & 34.2 & 16.9 & 37.0 & 7.0\\
Qwen-0.5B + LoRA & 7.6 & 8.2 & 45.9 & 34.8 & 5.6\\
\bottomrule\end{tabular}
\caption{Automatic error profile on the 500-pair test subset (\%). \emph{Exact}: output equals the reference. \emph{Kept}: dialect words (source words absent from the reference) left unchanged. \emph{Produced}: reference words absent from the source that the output contains. \emph{Novel}: output words in neither source nor reference. \emph{Neg.}: output and reference disagree on negation.}\label{tab:errors}
\end{table}

\begin{table*}[t]\centering\small
\begin{tabular}{@{}ll@{}}\toprule
\multicolumn{2}{@{}l}{\emph{Chittagong}}\\
Source & \bn{আর বর বইনোর আজিয়া মন ভালা নাই}\\
Reference & \bn{আমার বড় বোনের আজকে মন ভালো নেই}\\
3B 5-shot & \bn{এবং বর বইনোর জন্য মন ভালো নয়}\\
0.5B+LoRA & \bn{আমার বড় বোনের আজ মন ভালো নেই}\\
Dictionary & \bn{আমার খুব বোন আজ মন ভালো নেই}\\
BanglaT5 & \bn{আমার বর বোনের আজ মন ভালো নেই}\\
\midrule
\multicolumn{2}{@{}l}{\emph{Sylhet}}\\
Source & \bn{আমার বাসার দারওয়ান অনেক ভালা মানুষ}\\
Reference & \bn{আমার বাসার দারওয়ান অনেক ভালো মানুষ}\\
3B 5-shot & \bn{আমার বাসার দারোয়ান অনেক ভাল আমেরিকা}\\
0.5B+LoRA & \bn{আমার বাসার প্রধান অনেক ভালো মানুষ}\\
Dictionary & \bn{আমার বাসার দারওয়ান অনেক ভালো মানুষ}\\
BanglaT5 & \bn{আমার বাসার দারোয়ান অনেক ভালো মানুষ}\\
\midrule
\multicolumn{2}{@{}l}{\emph{Barishal}}\\
Source & \bn{মাইয়াডা বেইন্নাহালে ঘুমে গোনে উইঠ্যাই পড়তে বয়}\\
Reference & \bn{মেয়েটি সকালে ঘুম থেকে উঠেই পরতে বসে}\\
3B 5-shot & \bn{মেয়েটি বাইনহালে শুয়ে আছে এবং বয়ে যাওয়ার জন্য উঠে পড়তে বয়}\\
0.5B+LoRA & \bn{মেয়েটি সবার ঘুমিতে থাকতে বসে}\\
Dictionary & \bn{মেয়েটি সকালে ঘুম থেকে উইঠ্যাই পড়তে বসে}\\
BanglaT5 & \bn{মেয়েটি বিছানায় ঘুমিয়ে পড়ে}\\
\bottomrule\end{tabular}
\caption{Real test outputs, one sentence for three of the dialects (chosen automatically: the first subset sentence of 4--7 words on which the four systems all differ).}\label{tab:qual}
\end{table*}


\section{Discussion}
\paragraph{Why prompting fails.} The prompted 3B model treats unknown dialect words either as standard words,
which it keeps, or as noise, which it replaces with unrelated words. Five examples cannot teach a vocabulary
of thousands of word forms. Fine-tuning on 9,370 pairs teaches it, which is why even a 0.5B LLM gains
«ft_sub_gain» chrF over prompting.

\paragraph{Why simple systems win.} The task is close to word substitution: 15--18\% of the dictionary and
BanglaT5 outputs are exactly the reference, and the words that change are mostly a closed set of pronouns,
verb endings and frequent content words. A lookup table stores every observed replacement exactly, and
BanglaT5 is pretrained on Bangla and copies unchanged words reliably. A general LLM must instead learn the
replacements through low-rank adapters and override a prior that prefers paraphrase, which the error profile
shows as invented words. With «ft_short», LoRA on a 0.5B model does not close this gap; larger models and
longer training, as planned with QLoRA on Qwen2.5-7B, may.

\paragraph{Practical advice.} For normalising dialect text, a dictionary learned from parallel data, or a
fine-tuned BanglaT5 when out-of-domain text is expected, is a strong and cheap first choice. LLM-based systems
should be compared against these baselines, not only against copying or prompting. The copy-based distance
should be reported with any dialect result, because it predicts most of the variation between dialects and
between sentences.

\section{Limitations}
\paragraph{Compute.} All experiments ran on four CPU cores without a GPU. The planned main systems,
QLoRA-tuned Qwen2.5-7B-Instruct and Llama-3.1-8B-Instruct, and Llama-3.3-70B prompting, are implemented in the
repository (a Kaggle notebook runs them end to end) but are not reported here; the fine-tuned model is 0.5B and
the prompted model 3B. BanglaT5 had one epoch instead of eight, the fine-tuned LLM one seed, and the
leave-one-dialect-out study uses the dictionary only. Our conclusions about prompting therefore apply to
small open models; larger models may prompt better.

\paragraph{Contamination.} Vashantor has been public since 2023, so pretrained models may have seen it.
ONUBAD is not fully independent of Vashantor (323 shared standard sentences), so we use its novel part.

\paragraph{Data and metrics.} Vashantor sentences are short and conversational, and its test set covers a
limited set of topics; ONUBAD covers three dialects. chrF rewards surface overlap and can miss meaning errors
such as a changed negation; LaBSE and our negation check catch some of these, but we have no human ratings, so
``meaning preserved'' is measured only automatically.

\section{Conclusion}
We built a leakage-free benchmark from Vashantor and the novel part of ONUBAD and compared prompting,
fine-tuning and simple baselines for translating five Bangla dialects into standard Bangla. Prompting a 3B
instruction-tuned LLM is no better than copying the input, and LoRA fine-tuning of a 0.5B LLM improves on it by
«ft_sub_gain» chrF (RQ1), but a word-replacement dictionary and BanglaT5 remain clearly better, and BanglaT5
generalises best out of domain. Copy-based dialect distance predicts which dialects and which sentences are hard
for every system (RQ2), and about half of the dictionary knowledge needed for a dialect can be borrowed from the
other four (RQ3). The next steps are the GPU runs already implemented in the repository (QLoRA on Qwen2.5-7B and
Llama-3.1-8B, leave-one-dialect-out fine-tuning) and a human evaluation of meaning preservation.

{\raggedright\paragraph{Reproducibility.} \texttt{src/prepare\_data.py} and \texttt{src/prepare\_onubad.py} build the
splits; \texttt{src/baselines.py}, \texttt{src/dictionary\_experiments.py}, \texttt{src/generate.py},
\texttt{src/train\_qlora.py} and \texttt{src/train\_banglat5.py} produce the outputs; \texttt{src/evaluate.py}
and \texttt{src/auto\_error\_analysis.py} produce every table in \texttt{results/}; \texttt{scripts/paper\_tables.sh} and \texttt{report/build\_report.py} rebuild this paper.\par}

{\small
\begin{thebibliography}{11}
\bibitem{vashantor} F.~T.~J. Faria, M.~B. Moin, A.~A. Wase, M.~Ahmmed, M.~R. Sani, T.~Muhammad. Vashantor: A
large-scale multilingual benchmark dataset for automated translation of Bangla regional dialects to Bangla
language. arXiv:2311.11142, 2023.
\bibitem{banglat5} A.~Bhattacharjee et al. BanglaNLG and BanglaT5: Benchmarks and resources for evaluating
low-resource natural language generation in Bangla. Findings of EACL, 2023.
\bibitem{lora} E.~J. Hu et al. LoRA: Low-rank adaptation of large language models. ICLR, 2022.
\bibitem{qlora} T.~Dettmers, A.~Pagnoni, A.~Holtzman, L.~Zettlemoyer. QLoRA: Efficient finetuning of quantized
LLMs. NeurIPS, 2023.
\bibitem{chrf} M.~Popović. chrF: character n-gram F-score for automatic MT evaluation. WMT, 2015.
\bibitem{labse} F.~Feng, Y.~Yang, D.~Cer, N.~Arivazhagan, W.~Wang. Language-agnostic BERT sentence embedding.
ACL, 2022.
\bibitem{onubad} N.~Sultana, R.~Yasmin, B.~Mallik, M.~S. Uddin. ONUBAD: A comprehensive dataset for automated
conversion of Bangla regional dialects into standard Bengali dialect. Data in Brief, 58:111276, 2025.
\bibitem{jawad} M.~M. Jawad et al. Benchmarking large language models on Bangla dialect translation and
dialectal sentiment analysis. BLP-2025, pp.~322--337, 2025.
\bibitem{sami} K.~M.~J. Sami, D.~Sumit, A.~Hossain, F.~Sadeque. A comparative analysis of retrieval-augmented
generation techniques for Bengali standard-to-dialect machine translation using LLMs. arXiv:2512.14179, 2025.
\bibitem{qwen} Qwen Team. Qwen2.5 technical report. arXiv:2412.15115, 2024.
\end{thebibliography}}
\end{document}
